\chapter{Related Work}
\label{ch:related-work}

This chapter reviews the background that motivates combining structured knowledge graphs with explicit statistical dependence estimates:
biomedical knowledge graphs, graph neural networks for heterogeneous relational data, 
Hierarchical Correlation Reconstruction, and knowledge-aware recommendation with evaluation reliability.

\section{Knowledge Graphs in Biomedicine}
\label{sec:related-kg}

The term \emph{knowledge graph} (KG), though traceable to proposals from the
1970s, entered widespread use after Google's 2012 announcement of its
knowledge graph, built on resources such as DBpedia and Freebase
\cite{hogan2021knowledge}. Since then, the concept has been adopted far
beyond web search, as it provides an interlinked representation of
knowledge in a human- and machine-understandable
format~\cite{barrasa2023building}.

This growth has been particularly visible in biomedicine, where
KGs are used to integrate heterogeneous sources such as
drugs, genes, diseases, and mechanisms of action (the biochemical process by which a drug produces its effect) 
~\cite{maclean2021knowledge}.
In biomedical KGs, entities (nodes $V$) may represent drugs, diseases, genes,
proteins, mechanisms, patient characteristics, or clinical outcomes, while
edges ($E$) describe typed relationships (e.g.,
\emph{drug--treats--disease}, \emph{gene--associates--disease}).
Such graphs are naturally heterogeneous and admit the form~\cite{hu2020hgt}
\begin{equation}
\label{eq:hetero-graph}
\begin{gathered}
G = \bigl(V, E, \tau, \phi\bigr), \\[0.3em]
{\scriptstyle
\tau \colon V \longrightarrow \mathcal{T}_V \text{ and }
\phi \colon E \longrightarrow \mathcal{T}_E
}
\end{gathered}
\end{equation}

Within this domain, knowledge graphs support a wide range of tasks,
spanning drug-drug interaction prediction, drug development, 
and clinical decision support~\cite{cui2025hkgreview}.
Among these, clinical decision support centres on
patient-level data as a core, rather than drug- or trial-related data.
Graph-based approaches within this patient-centric
setting have been applied to different tasks such as clinical prediction tasks including
in-hospital mortality and readmission~\cite{jiang2024graphcare}, survival
analysis~\cite{tan2023gnnsurv}, and multimodal risk prediction from
electronic health records (EHR)~\cite{zitnik2024graphai}.

At this point it is important to note that such clinical prediction models
could pursue different aims, such as risk prediction or treatment effect estimation. 
Risk models
predict the probability of an outcome under the patient's observed
factual course of care, whereas treatment effect estimation asks a
counterfactual question: how would the outcome differ under an
alternative treatment decision~\cite{hoogland2021tutorial}. This distinction motivates
 the design of our knowledge graph, which is structured to reflect 
 intended causal relations between entities and patient trajectories in 
 relation to specific treatments (drugs) and patient characteristics.
The decision was motivated by Toonsi et al.~\cite{toonsi2025causal} who propose
causal
knowledge graphs as a means of addressing several shortcomings of the
standard modelling pipeline simultaneously. 
As KG's edges are
intended to encode causal, rather than purely associative, relations,
causal knowledge graphs are more robust to distribution
shift between clinics and better suited to reasoning under incomplete
documentation, while also allowing prior medical knowledge to support
downstream prediction tasks~\cite{cui2025hkgreview}. 

Such graphs could be dynamically constructed from empirical data.
Standard pipelines for constructing KG include a dedicated
step for completing missing relations between entities using graph
databases or link prediction models, typically applied to graphs
derived from literature and curated sources~\cite{cui2025hkgreview}. 

%This
%thesis revisits the same construction step applying link prediction 
%to reconstruct relations within graphs
%built from empirical data.

Taken together, these observations motivate two complementary tasks
studied in this thesis: reconstructing
knowledge graphs via empirical patient-level data, and predicting clinical outcomes
on the resulting structure via node classification. The
tasks are addressed by GNNs described in the following section.

\subsection{Graph Neural Networks}
\label{subsec:related-gnn}

Graph neural networks (GNNs) learn node representations by aggregating
information from neighbouring nodes through message passing
\cite{hamilton2017graphsage}.
After $K$ message-passing layers, the representation of a node can
incorporate information from its $K$-hop neighbourhood.
During training, every node is represented not only by its own features, as
in traditional neural networks, but also by the features of its neighbours.
Unlike earlier graph-embedding methods with fixed representations, GNNs
build node embeddings inductively through aggregation, so inference does not
require all nodes to be present during training.

Formally, the representation $h_v^{(K)}$ of node $v$ depends on the features
of all nodes within its $K$-hop neighbourhood.
At layer $k$,
\begin{equation}
\label{eq:related-mp}
\begin{gathered}
h_v^{(k)} = \text{UPDATE}\Bigl(h_v^{(k-1)},\ \text{AGGREGATE}\bigl(\{h_u^{(k-1)} : u \in \mathcal{N}(v)\}\bigr)\Bigr), \\[0.3em]
{\scriptstyle
\mathcal{N}(v) \text{ --- neighbourhood of node } v
}
\end{gathered}
\end{equation}
so that information from a node $u$ at graph distance $d$ from $v$ only
reaches $v$ once $K \geq d$.
Different GNN architectures mainly differ in how neighbouring information
and, on heterogeneous graphs, relation types are transformed and aggregated.
Several architectures relevant to this thesis are reviewed below.

\subsubsection{GraphSAGE}

GraphSAGE~\cite{hamilton2017graphsage} was introduced for inductive and
scalable representation learning on large graphs.
It instantiates the general message-passing scheme using one of three
aggregation functions---mean, LSTM, or pooling---so that the same learned
functions can later be applied to unseen nodes.
GraphSAGE concatenates a node's own representation with the aggregated
neighbour representation and L2-normalises the result after each layer.
In order to improve scalability, it samples a fixed number of neighbours per node.

\subsubsection{Relational Graph Convolutional Network}

Relational Graph Convolutional Networks (R-GCNs)
\cite{schlichtkrull2018rgcn} extend graph convolutional networks to
multi-relational graphs, in which edges represent different types of
relationships.
Unlike a standard GCN, which applies the same transformation to messages
from all neighbours, R-GCN applies a relation-specific transformation.

For a node $v$, the representation at layer $l+1$ is computed as
\begin{equation}
\label{eq:rgcn}
\begin{gathered}
h_v^{(l+1)} =
\sigma \left(
\sum_{r \in \mathcal{R}}
\sum_{u \in \mathcal{N}_r(v)}
\frac{1}{c_{v,r}}
W_r^{(l)} h_u^{(l)}
+
W_0^{(l)} h_v^{(l)}
\right), \\[0.4em]
{\scriptstyle
\mathcal{R}\text{ --- relation types},\ 
\mathcal{N}_r(v)\text{ --- neighbours of }v\text{ under }r,\ 
c_{v,r}\text{ --- normalisation}
} \\[0.15em]
{\scriptstyle
W_r^{(l)}\text{ --- relation-specific weights},\ 
W_0^{(l)}h_v^{(l)}\text{ --- self-loop},\ 
\sigma\text{ --- activation}
}
\end{gathered}
\end{equation}

When the number of relation types is large, assigning a separate full weight
matrix to each relation can lead to overfitting.
R-GCN therefore introduces basis decomposition, in which each
relation-specific matrix is expressed as a linear combination of $B$ shared
basis matrices:
\begin{equation}
\label{eq:rgcn-basis}
\begin{gathered}
W_r^{(l)} =
\sum_{b=1}^{B} a_{rb}^{(l)} V_b^{(l)}, \\[0.4em]
{\scriptstyle
V_b^{(l)}\in\mathbb{R}^{d_{\mathrm{out}}\times d_{\mathrm{in}}}\text{ --- shared basis},\ 
a_{rb}^{(l)}\text{ --- coefficient for relation }r
}
\end{gathered}
\end{equation}

\subsubsection{Transformer-based message passing}

%Attention-based GNNs learn the importance of neighbouring nodes rather than
%using a fixed aggregation rule \cite{velickovic2018graphattention}.
Transformer-based graph neural networks adapt the self-attention mechanism
from the Transformer architecture of Vaswani et al.~\cite{vaswani2017attention} to
graph-structured data.
For each node $i$, its input representation $\mathbf{h}_i$ is projected into
query, key, and value vectors:
\begin{equation}
\label{eq:transformer-qkv}
\begin{gathered}
\mathbf{q}_i = \mathbf{W}_{Q}\mathbf{h}_i, \qquad
\mathbf{k}_i = \mathbf{W}_{K}\mathbf{h}_i, \qquad
\mathbf{v}_i = \mathbf{W}_{V}\mathbf{h}_i, \\[0.4em]
{\scriptstyle
\mathbf{W}_{Q},\;\mathbf{W}_{K},\;\mathbf{W}_{V}\text{ --- learnable projection matrices}
}
\end{gathered}
\end{equation}
For an edge from source node $j$ to target node $i$, the attention
coefficient is
\begin{equation}
\begin{gathered}
e_{ij} =
\frac{\mathbf{q}_{i}^{\top}\mathbf{k}_{j}}{\sqrt{d_k}}, \\[0.4em]
{\scriptstyle
d_k\text{ --- dimensionality of the key vectors}
}
\end{gathered}
\end{equation}
The attention coefficients are normalised across all neighbours of the
target node using the softmax function:
\begin{equation}
\label{eq:transformer-softmax}
\begin{gathered}
\alpha_{ij} =
\frac{\exp(e_{ij})}
{\sum_{r \in \mathcal{N}(i)} \exp(e_{ir})}, \\[0.4em]
{\scriptstyle
\mathcal{N}(i)\text{ --- neighbourhood of target node }i
}
\end{gathered}
\end{equation}
The node's representation is updated by aggregating value projections of
its neighbours, weighted by the learned attention coefficients:
\begin{equation}
\begin{gathered}
\mathbf{h}'_i =
\sum_{j \in \mathcal{N}(i)}
\alpha_{ij}\mathbf{v}_{j}, \\[0.4em]
{\scriptstyle
\alpha_{ij}\text{ --- attention weight},\ 
\mathbf{v}_{j}\text{ --- value projection of neighbour }j
}
\end{gathered}
\end{equation}

To capture multiple complementary interaction patterns, graph Transformers
commonly use multi-head attention.
The output of head $m$ for node $i$ is:
\begin{equation}
\begin{gathered}
\mathbf{h}_i^{\prime(m)} =
\sum_{j \in \mathcal{N}(i)}
\alpha_{ij}^{(m)}\mathbf{v}_{j}^{(m)}, \\[0.4em]
{\scriptstyle
m\in\{1,\ldots,H\},\ 
\mathbf{W}_{Q}^{(m)},\;\mathbf{W}_{K}^{(m)},\;\mathbf{W}_{V}^{(m)}\text{ --- independent projections per head}
}
\end{gathered}
\end{equation}
The outputs of all $H$ heads are combined by concatenation or averaging.

All architectures mentioned above are homogeneous from nature, but can be applied to heterogeneous
graphs through an appropriate configuration of the \texttt{HeteroConv} layer in
PyTorch Geometric (PyG). \texttt{HeteroConv} is a generic wrapper that applies
a separately specified message-passing operator to each edge type, i.e., each
meta-relation between a source node type and a target node type~\cite{wojcik2026grafowe}.
This mechanism is used in Task B which allowed to experiment with different graph architectures 
on joined setting.


\subsubsection{Heterogeneous Graph Transformer}

In response to the demand for preserving heterogeneous structure during modelling, Hu et al.~\cite{hu2020hgt} 
proposed the Heterogeneous Graph Transformer (HGT).
For attention-based message passing on heterogeneous graphs, the
Heterogeneous Graph Transformer (HGT) makes attention scores and message
transformations dependent on the types of source and target nodes and on the
edge relation.
Unlike the \texttt{HeteroConv}-wrapped TransformerConv described above,
which instantiates a fully independent set of attention parameters for
each edge type and aggregates their outputs only after each layer has
been computed separately, HGT shares a single underlying attention
mechanism across all edge and node types, using type-specific projection
matrices. This allows HGT to learn
interactions across types directly, at a lower parameter cost as the
number of types grows~\cite{hu2020hgt}.

The HGT layer decomposes message passing
into three components: heterogeneous mutual attention, heterogeneous
message passing, and target-specific aggregation. Each component uses a
separate set of weight matrices for every combination of source node
type, edge type, and target node type -- referred to as a
\emph{meta relation} $(\tau_s, \phi(e), \tau_t)$.
%For an edge from source node $s$ to target node $t$, HGT represents its
%semantics as a meta-relation $(\tau_s, \phi(e), \tau_t)$, where $\tau_s$
%and $\tau_t$ denote the source and target node types and $\phi(e)$ denotes
%the edge type.
So that HGT extends the query, key, and value projections introduced in
Eq.~\eqref{eq:transformer-qkv} by conditioning each weight matrix on node
type: $\mathbf{W}_{Q}^{\tau_t}$, $\mathbf{W}_{K}^{\tau_s}$, and
$\mathbf{W}_{M}^{\tau_s}$ replace the shared $\mathbf{W}_{Q}$,
$\mathbf{W}_{K}$, and $\mathbf{W}_{V}$, so that
\begin{equation}
\mathbf{q}_t = \mathbf{W}_{Q}^{\tau_t}\mathbf{h}_t, \qquad
\mathbf{k}_s = \mathbf{W}_{K}^{\tau_s}\mathbf{h}_s, \qquad
\mathbf{m}_s = \mathbf{W}_{M}^{\tau_s}\mathbf{h}_s.
\end{equation}
The attention coefficient additionally depends on the edge type through a
relation-specific matrix $\mathbf{W}^{\phi(e)}_{\text{ATT}}$:
\begin{equation}
e_{ts} =
\frac{\left(\mathbf{q}_t^{\top}\mathbf{W}^{\phi(e)}_{\text{ATT}}\mathbf{k}_s\right)}
{\sqrt{d_k}},
\qquad
\alpha_{ts} = \frac{\exp(e_{ts})}{\sum_{r \in \mathcal{N}(t)} \exp(e_{tr})}.
\end{equation}
The message is similarly transformed by a relation-specific matrix
$\mathbf{W}^{\phi(e)}_{\text{MSG}}$ and aggregated across all neighbours
of $t$:
\begin{equation}
\label{eq:hgt-aggregation}
\mathbf{h}'_t =
\sum_{s \in \mathcal{N}(t)}
\alpha_{ts}\left(\mathbf{m}_s\mathbf{W}^{\phi(e)}_{\text{MSG}}\right)~\cite{hu2020hgt}.
\end{equation}
Relation-specific transformations affect both the attention
score and the transferred message, so that HGT assigns different
importance to neighbours based not only on their features, but also on
the semantic type of their connection to the target node.

As $\mathbf{h}'_t$ combines information from source nodes of
potentially different types and feature distributions, it is mapped back
into the target node's type-specific space via a target-type-dependent
projection, followed by a non-linear activation and a residual
connection to the previous layer:
\begin{equation}
\mathbf{h}''_t =
\mathbf{A}\text{-}\mathrm{Linear}_{\tau_t}
\left(\sigma\left(\mathbf{h}'_t\right)\right)
+ \mathbf{h}_t^{(l-1)}~\cite{hu2020hgt}.
\end{equation}

This capacity to differentiate neighbours by both node type and
relation semantics is particularly relevant to the knowledge graphs
built for the purpose of this thesis, where clinically distinct relation types --
such as drug--drug interactions versus drug--outcome relations -- 
 carry different predictive weight. 
HGT is the primary structural encoder in Task~A and proxy to Task~A, because typed
pharmacotherapy and collaborative-knowledge-graph relations require type-aware message
passing. 



% DropEdge (Rong et al., 2020) randomly removes a fraction of edges independently in every layer during training, slowing the rate of information mixing and acting as a structural regularizer.

% Jumping Knowledge (Xu et al., 2018) aggregates the outputs of all layers (via concatenation, cat, or elementwise max) rather than only the last one, letting a node fall back on a shallower, less-smoothed representation when the deepest layer has become uninformative.

% Weisfeiler-Lehman --- two nodes with identical initial features and structurally equivalent neighborhoods 
% cannot be assigned distinct representations by any number of convolutional layers.

% leave-one-out

\section{Statistical Dependence and Hierarchical Correlation Reconstruction}
\label{sec:related-hcr}
\label{subsec:hcr_def}

Empirical comparisons between graph-based and tabular models report
mixed results, with gains from graph structure often modest and
architecture-dependent~\cite{tabgraphs2024,zhu2025stroketh}. Rather than
resolving this solely through architecture search, this thesis follows
the Wide\,\&\,Deep paradigm~\cite{cheng2016wideanddeep,gao2021wdgnn}, in
which a linear or statistically-grounded branch complements a deep,
graph-based branch. The wide branch memorises direct, low-order
statistical dependencies that message passing may dilute through
neighbourhood aggregation, while the deep branch generalises through
learned relational representations.  
In this thesis, the wide branch is introduced using
Hierarchical Correlation Reconstruction (HCR) described in this section.

HCR, proposed by Jaros\l{}aw Duda, represents probability distributions and statistical dependencies through coefficients of an orthonormal basis expansion~\cite{duda2018hcr,duda2018hcr-ts,duda2020hcr-income,duda2023hcr-wolfram,duda2025hcr-mi}.
Lower-order coefficients describe broad dependence patterns; higher-order coefficients capture more complex departures from independence.

For variables $X$ and $Y$ with orthonormal bases $\{\phi_j^{(X)}\}$ and $\{\phi_k^{(Y)}\}$, empirical cross-dependence coefficients are
\begin{equation}
a_{jk}
=
\frac{1}{n}
\sum_{i=1}^{n}
\phi_j^{(X)}(X_i)\,
\phi_k^{(Y)}(Y_i).
\end{equation}
For a binary variable $B\in\{0,1\}$ with prevalence $p_B=P(B=1)$,
\begin{equation}
\phi_1(B)=\frac{B-p_B}{\sqrt{p_B(1-p_B)}},
\end{equation}
so the lowest-order coefficient $a_{11}$ is a standardized first-order dependence descriptor.
The same principle extends to other variable types via suitable marginal bases.

HCR has demonstrated practical value in domains requiring reliable,
interpretable estimates of conditional distributions instead of single
point predictions. In virtual screening for drug discovery, HCR-based
prediction of ADMET property distributions enables inexpensive
selection of compounds with properties nearly certain to fall within a
desired range, while automatically flagging predictions with high
uncertainty for manual review~\cite{szarek2022hcrmolecular}. In income
credibility assessment using Polish household budget survey data, HCR
produced more accurate proxies for low-income households than
conventional tax-record matching, showing its ability to model the
conditional distribution directly instead of relying on a single
regression estimate~\cite{duda2020hcr-income,brzezinski2022polishinequality}.

Based on those results, in this work HCR is used as a statistical representation of empirical and 
associative dependence to complement the GNN representation.
Causal direction and
mechanistic interpretation are instead supplied by the directed graph
schema itself built for the purpose of this thesis.
In general the term HCR denotes the orthonormal-basis coefficient construction, whereas this thesis
uses HCR-inspired representations -- larger task-specific vectors that 
combine those coefficients with additional descriptors.
In Task~A they describe the local dependency motif around a candidate edge, while in Task~B they supply 
an endpoint-neighbourhood evidence channel alongside the GNN representation.
The common principle is combining learned graph representations with explicit statistical 
dependence estimated from observations.
